\section{A constructive phase-rounding certificate}
Let $F$ be a complex Hadamard matrix of order $n$, with entries of
modulus one and $FF^*=nI$. Set $M=F^{\circ4}$ and
\[
 D_4(F)=n^4-\|MM^*\|_F^2.
\]
The nonnegative deficit $D_4$ is invariant under row/column phases
and permutations. The previous pass proved an obstruction when
$\rho_4=\|M\|_{\rm op}/n$ is bounded away from one. Here we prove a
positive, dimension-dependent rounding statement.

\begin{theorem}[Finite dephasing and exact quaternary rounding]
\label{thm:phase-rounding}
If $D_4(F)<n$, or $D_4(F)<2n$ when $n$ is even,
row and column phases convert $F$ into
a matrix whose entrywise nearest-fourth-root rounding is a
quaternary Hadamard matrix of order $n$.
Only $n^2$ explicit dephasings need be tested.
In particular it suffices that
\[
 1-\rho_4(F)^4<
 \begin{cases}2/n^3,& n\text{ even},\\1/n^3,&n\text{ odd}.\end{cases}
\]
This sufficient deficit is of order $n^{-3}$ in the last inequality;
no such estimate is established at every even order.
\end{theorem}
\begin{proof}
For each pivot $(k,\ell)$ define the dephased matrix
\[
 F^{k\ell}_{ij}
 =F_{ij}\overline{F_{kj}}\,
       \overline{F_{i\ell}}F_{k\ell}.
\]
These are row and column phase operations, so each matrix is still
Hadamard. Its entrywise fourth power is the corresponding four-cycle
cross ratio of $M$. Direct expansion and summation give the exact identity
\[
 \sum_{k,\ell}\|(F^{k\ell})^{\circ4}-J_n\|_F^2=2D_4(F).
\]
Indeed the sum of the real parts of those cross ratios is
$\sum_{i,k}|\sum_j M_{ij}\overline{M_{kj}}|^2=\|MM^*\|_F^2$.
Some pivot therefore has fourth-power squared error at most
$2D_4/n^2$.

For a unit complex number $z$, let $q\in\{1,-1,i,-i\}$ be nearest.
Write $z/q=e^{it}$ with $|t|\le\pi/4$. Then
\[
 |z-q|^2=4\sin^2(t/2)\le
 \tfrac14\,4\sin^2(2t)=\tfrac14|z^4-1|^2.
\]
For example $\sin(2t)=4\sin(t/2)\cos(t/2)\cos t$, whose absolute
ratio to $\sin(t/2)$ is at least two on this interval.
Thus the rounded matrix $Q$ at that pivot has
$\eta=\|F^{k\ell}-Q\|_F^2\le D_4/(2n^2)$.
Write $e_i$ for the difference between row $i$ of $Q$ and the
corresponding row of $F^{k\ell}$. For distinct rows, orthogonality
and the fact that both matrices have row norm $\sqrt n$ give
\[
 |(QQ^*)_{ij}|
 \le\sqrt n(\|e_i\|+\|e_j\|)
 \le\sqrt{2n\eta}.
\]
Every off-diagonal Gram entry is a Gaussian integer. A nonzero one
has modulus at least one, so $\eta<1/(2n)$ suffices to force zero.
For even $n$, its real and imaginary integer parts have even sum:
it is a sum of an even number of fourth roots. A nonzero such entry
has modulus at least $\sqrt2$, so $\eta<1/n$ suffices.
The diagonal Gram entries are already $n$. This proves exact
quaternary orthogonality under the stated deficit bounds. Finally
$\|MM^*\|_F^2\ge\|M\|_{\rm op}^4=n^4\rho_4^4$ proves the last claim.
\end{proof}

The proof is an algorithm: dephase at all pivots, choose one minimizing
the fourth-power error, round, and check the Gaussian-integer Gram.
Alternatively try every pivot and stop at an exact Gram certificate.
No target Hadamard matrix is needed as an input. Computing all pivot
errors naively uses $O(n^4)$ operations. Computing row inner products
of $M$ first permits the error scores to be obtained in $O(n^3)$
operations: the score at $(k,\ell)$ is
\[
 2n^2-2\operatorname{Re}\sum_i
   \overline{M_{i\ell}}M_{k\ell}
   \sum_j M_{ij}\overline{M_{kj}}.
\]
Rounding and checking one selected pivot use $O(n^3)$ additional
operations. Arithmetic complexity here counts exact scalar operations;
it is not a bit-complexity bound for arbitrary real inputs.

\begin{corollary}[Separation in the absence of a quaternary matrix]
If $\mathrm{BH}(n,4)$ does not exist, every complex Hadamard matrix
of that order satisfies
\[
 D_4(F)\ge
 \begin{cases}2n,&n\text{ even},\\n,&n\text{ odd},\end{cases}
 \qquad
 1-\rho_4(F)^4\ge
 \begin{cases}2/n^3,&n\text{ even},\\1/n^3,&n\text{ odd}.\end{cases}
\]
\end{corollary}
This is a contrapositive separation theorem, not a construction of a
point below the threshold. At odd $n>1$ quaternary existence is
impossible: a vanishing sum of an odd number of fourth roots has
unequal parity in its real and imaginary counts. At even $n$,
an actual threshold-crossing construction remains the central task.
The classical fourth-cycle invariants and Butson realification are
background; see Tadej--\.Zyczkowski,
\url{https://arxiv.org/abs/quant-ph/0512154}.
No bibliographic novelty claim is made for the rounding estimate.

\subsection{A certificate that does not assume exact orthogonality}
The same algorithm also accepts approximate inputs.
Let $F$ now be any $n$-square unit-modulus matrix, define $D_4$ as
before, and put $\beta=\max_{i\ne j}|(FF^*)_{ij}|$.
Define $\gamma_n=\sqrt2$ for even $n$ and $\gamma_n=1$ otherwise.
\begin{proposition}[Two-objective certificate]
If
\[
 \beta+\sqrt{D_4(F)/n}<\gamma_n,
\]
one of the pivot dephasings rounds to an exact $\mathrm{BH}(n,4)$.
\end{proposition}
\begin{proof}
The averaging identity and nearest-root estimate only use unit modulus.
Dephasing preserves the magnitudes of the original Gram entries.
For the selected pivot and its rounded $Q$, the same two-row
estimate, now retaining the original correlation, gives
\[
 |(QQ^*)_{ij}|\le\beta+\sqrt{2n\eta}
 \le\beta+\sqrt{D_4/n}<\gamma_n.
\]
Gaussian-integrality and parity again force zero.
\end{proof}
Thus a continuous search could minimize correlation and phase dispersion
jointly, then obtain a discrete exact certificate on crossing this
explicit boundary. This is not a convergence theorem. In particular a
real sign matrix already has $D_4=0$, but its even integer correlations
must already vanish to satisfy this sufficient condition. The certificate
does not turn arbitrary sign-matrix local minima into Hadamard matrices.
